\section{Marco Teórico}

\subsection{EPS e IPS}

Puede llegar a existir la confusión entre estos dos términos los cuales provienen de la Ley 100 la cual reglamenta el sistema de salud colombiano y provee las herramientas para la financiación y sostenibilidad de su estructura procurando garantizar una atención igualitaria para todas las personas.

En primer lugar, las EPS, Entidades Promotoras de Salud, son las responsables de la afiliación y el registro de los afiliados y del recaudo de sus cotizaciones. Como dice artículo publicado por la Alcaldía de Bogotá:

\begin{quote}
	Su función básica es organizar y garantizar la prestación del Plan Obligatorio de Salud (POS) y hacer los giros respectivos al Fondo de Solidaridad y Garantía que es donde se administran los recursos del Sistema de Seguridad Social en Salud. Así, todas las personas se afilian a estas y quedan amparados en su intermediación para acceder a los servicios médicos. \cite{arias2021}
\end{quote}

Por otro lado, las IPS, Instituciones Prestadoras de Servicios de Salud, corresponden a todas las entidades, asociaciones y/o personas bien sean públicas, privadas o con economía mixta, que están autorizadas para prestar de forma parcial y/o total los procedimientos que se demanden para cumplir el Plan Obligatorio de Salud (POS). Dentro de este grupo es que entran los hospitales, clínicas y distintos centros de salud.

\subsection{Deuda de las EPS}

El año pasado, mediante un boletín de prensa, el número 093-2025, del día 15 de Julio, emitió el último informe (hasta ese momento) por parte de la Contraloría General de la República sobre la situación financiera del sistema de salud que revela 29 Entidades Promotoras de Salud (EPS) acumulan una deuda que asciende a los \$32,9 billones de pesos por servicios prestados por clínicas, hospitales, laboratorios, operadores farmacéuticos y otros actores del sistema.

Esta cifra, como bien se da en el informe dado por la Contraloría, representó un aumento de \$7.9 billones comparado al monto reportado en 2023, lo cual claramente refleja un desorden a nivel económico desde la parte administrativa en el que se encuentra todo el sistema de salud, que como expresan en el reporte, es ``debido al actual modelo de aseguramiento basado en la intermediación.'' \cite{minsalud2025}

Como consecuencia a ello, la deuda en primera instancia afecta gravemente la estabilidad de los prestadores y proveedores de servicios, pero se genera una cadena llegando a afectar también de sobremanera a los usuarios, en donde se llega a violar el derecho fundamental a la salud en escenarios como barrera en el acceso con los largos retrasos para citas y cirugías, el desabastecimiento de insumos médicos y tecnologías necesarias y llegando a declararlo ya como emergencia humanitaria (por parte de AFIDRO - Asociación de Laboratorios Farmacéuticos de Investigación).

Es válido mencionar que, dentro de este boletín, lo dicho es que la culpa no recae en el gobierno, sino que va directamente a las EPS, de sus decisiones, omisiones y su mala organización administrativa.

\subsection{Reconocimiento Óptico de Caracteres (OCR)}

El Reconocimiento Óptico de Caracteres (\textit{Optical Character Recognition}, OCR) es un conjunto de técnicas que permite \textbf{extraer texto desde imágenes}. En el contexto de \textit{MediFinder}, el OCR es especialmente relevante porque tanto las \textbf{cajas de medicamentos} como la \textbf{parte trasera del blíster} suelen incluir texto con información crítica (nombre, concentración, forma farmacéutica, laboratorio y en algunos casos registro).

En términos generales, un sistema de OCR aplicado a fotografías tomadas por usuarios suele incluir:
\begin{itemize}
	\item \textbf{Preprocesamiento:} corrección de inclinación/rotación, aumento de contraste, reducción de ruido y recorte de regiones con texto.
	\item \textbf{Detección de texto:} localizar regiones donde probablemente hay texto (por ejemplo, líneas o bloques).
	\item \textbf{Reconocimiento:} convertir píxeles a caracteres para obtener una cadena de texto.
	\item \textbf{Postprocesamiento:} limpieza y normalización del texto (tildes, mayúsculas, unidades como mg/mL, abreviaturas) y corrección de errores típicos (por ejemplo, confusiones entre ``O'' y ``0'').
\end{itemize}

En un MVP orientado al usuario, es una buena práctica incorporar una etapa de \textbf{verificación humana} (el usuario corrige el texto detectado) para disminuir los efectos de errores de OCR causados por baja iluminación, reflejos (comunes en blíster) o desenfoque.

\subsection{Redes Neuronales Convolucionales (CNN)}

Para poder definir las redes neuronales convolucionales, en primer lugar, toca dar paso a la definición y el funcionamiento de las redes neuronales en sí.

IBM, dentro de su página web, brinda la definición al respecto de este tema:

\begin{quote}
	Las redes neuronales son un subconjunto del machine learning y están en el corazón de los algoritmos de deep learning. Están compuestas por capas de nodos, que incluyen una capa de entrada, una o más capas ocultas y una capa de output. Cada nodo está conectado a otro y tiene un peso y un umbral asociados. \cite{ibm2025cnn}
\end{quote}

Es válido hacer la mención de la definición general de las redes neuronales debido a que existen varios tipos usados para distintos casos y para distintos tipos de datos. Por ejemplo, las redes neuronales recurrentes se suelen usar para el procesamiento del lenguaje natural y el reconocimiento de voz. Y están las redes neuronales convolucionales, que son utilizadas principalmente para tareas de clasificación y computer vision.

El funcionamiento de este tipo de redes neuronales es destacable en el área de imágenes, voz o audio, y se componen de tres tipos de capas (principalmente): capa convolucional, capa de agrupación y la capa totalmente conectada.

En \textit{MediFinder}, la CNN se contempla principalmente como un \textbf{mecanismo de respaldo} (fallback) en los casos donde el OCR falle (por ejemplo, texto ilegible por reflejos en el blíster). En ese escenario, la CNN puede extraer un \textbf{embedding visual} del empaque para recuperar candidatos cercanos por similitud y luego cruzarlos con metadata del catálogo.

\subsection{K-Nearest Neighbors (KNN)}

Por otro lado, el enfoque del proyecto no se basa únicamente en extraer texto (OCR), sino también en recomendar alternativas. Para esto, se usa el algoritmo \textit{k-nearest neighbors} (KNN) como método de \textbf{búsqueda por similitud} dentro de un catálogo de medicamentos, aprovechando atributos como principio activo principal, concentración, forma farmacéutica o representaciones vectoriales (embeddings).

En su definición, el algoritmo de k-nearest neighbors (KNN) es un clasificador de aprendizaje supervisado no paramétrico, que emplea la proximidad para realizar clasificaciones o predicciones sobre la agrupación de un punto de datos individual. Hoy en día, es uno de los clasificadores de clasificación y regresión más sonados además de sencillos de implementar.

Ante la explicación precisa del funcionamiento del supuesto de que hay puntos similares cercanos, dentro de la página de IBM:

\begin{quote}
	En los problemas de clasificación, la etiqueta de clase se asigna por mayoría de votos, es decir, se emplea la etiqueta que se representa con más frecuencia en torno a un punto de datos determinado. Aunque técnicamente se considera ``votación por pluralidad'', el término ``votación por mayoría'' se emplea más comúnmente en las publicaciones al respecto. La distinción entre estas terminologías es que la ``votación por mayoría'' requiere técnicamente una mayoría superior al 50\%, lo que funciona sobre todo cuando sólo hay dos categorías. Cuando hay varias clases, por ejemplo, cuatro categorías, usted no requiere necesariamente el 50\% de los votos para llegar a una conclusión sobre una clase, ya que podría asignar una etiqueta de clase con un voto superior al 25\%. \cite{ibm2025knn}
\end{quote}

En \textit{MediFinder}, KNN se utiliza para \textbf{rankear alternativas} y proponer medicamentos similares, priorizando aquellos que comparten el \textbf{principio activo principal} (misma molécula). Como segundo nivel informativo (cuando el catálogo no ofrece suficientes equivalentes), se pueden mostrar alternativas de la misma \textbf{clase terapéutica} con advertencias, ya que no necesariamente son intercambiables.

\pagebreak